% How OCaml tells an unboxed int from a pointer: the low bit.
% Two machine words. Used in M10-L02.
\documentclass[tikz,border=4pt]{standalone}
\usepackage{tikz}
\input{_m10-styles.texinput}
\begin{document}
\begin{tikzpicture}[m10, x=1mm, y=-1mm, font=\large]
  % ---- immediate integer ----
  \node[note, anchor=west, font=\Large\bfseries] at (0,-4) {Immediate: the int 42};
  \node[heapcell, minimum width=85mm, minimum height=12mm, font=\large] at (42,8)
    {63-bit value \(= 42\)};
  \node[heapcell, minimum width=12mm, minimum height=12mm, fill=heapedge!20,
        font=\Large\bfseries] at (91,8) {1};
  \node[goodnote, anchor=west, font=\large] at (100,8) {low bit 1: it IS the value};

  % ---- pointer ----
  \node[note, anchor=west, font=\Large\bfseries] at (0,28) {Pointer: a boxed value};
  \node[stackcell, minimum width=85mm, minimum height=12mm, font=\large] at (42,40)
    {address of a heap block};
  \node[stackcell, minimum width=12mm, minimum height=12mm, fill=stackedge!20,
        font=\Large\bfseries] at (91,40) {0};
  \node[note, anchor=west, font=\large] at (100,40) {low bit 0: follow it};

  \node[note, anchor=north west, font=\large\itshape] at (0,56)
    {One bit-test, no type tags in memory. The cost: ints are 63-bit.};
\end{tikzpicture}
\end{document}
